\chapter{Life sciences}
\label{ch:life}

\section{The one sector that already owns something}

Everywhere else in Part III the problem is building a capability. Here the
capability exists, is forty years old, has products, has a reputation, and has
never been organised to sell anything.

Cuban biotechnology began in the early 1980s as a deliberate state decision to
build a research capacity in a poor country, and it produced, against every
reasonable expectation, a genuine one. The institutions --- the Centre for
Genetic Engineering and Biotechnology, the Finlay Institute, the Centre of
Molecular Immunology, and the rest, consolidated since 2012 into the
BioCubaFarma holding --- have delivered products that were not derivative.

The meningococcal B vaccine developed at the Finlay Institute in the 1980s was
the first of its kind anywhere. The \emph{Haemophilus influenzae} type b vaccine
licensed in 2004 used a fully synthetic antigen, the first synthetic-antigen
vaccine in routine human use. Heberprot-P, a recombinant epidermal growth factor
injected into diabetic foot ulcers, addresses a condition that causes a large
number of the world's amputations and has no close equivalent elsewhere.
CIMAvax-EGF, a therapeutic lung cancer vaccine, has been sufficiently interesting
to American oncologists that a United States cancer centre pursued trials of it
through the licensing exceptions. And in 2021 the sector developed, tested,
manufactured and deployed its own COVID-19 vaccines, vaccinating over ninety per
cent of the population --- the only country in Latin America to do so, and it did
it while unable to reliably import syringes.\unv{}

Why did this work when nothing else did? Three reasons, and they are worth
naming because the reform must not destroy them.

It was a \emph{closed loop}. Research, clinical development, manufacture and
deployment sat inside one system with one customer, the national health service,
so a product went from bench to population without an intervening market and
without a reimbursement negotiation. That is a genuine structural advantage for
developing products aimed at a whole population.

It had \emph{patient capital}. The state funded it for a decade before results,
which no shareholder would have done, and did not close it during the Special
Period.

And it had \emph{no profit constraint on target selection}. It worked on
meningitis B, on diabetic foot, on lung cancer in smokers --- conditions of
interest to Cuba and to poor countries generally rather than to the markets that
set global pharmaceutical priorities.

\section{Why it earns almost nothing}

And yet. Cuban pharmaceutical and biotechnology exports are, by any published
measure, a small fraction of what a portfolio like this should be worth --- of
the order of a few hundred million dollars a year, against a global market in the
trillions and against Cuban medical \emph{services} exports that were at their
peak twenty or thirty times larger.\unv{}

The reasons are not scientific. They are regulatory, contractual and financial,
and each is fixable.

\emph{No regulatory conversion.} Almost nothing in the portfolio holds approval
from the United States Food and Drug Administration or the European Medicines
Agency, and much of the clinical evidence was not generated to the standards
those agencies require --- not because the trials were dishonest, but because
they were designed for Cuban registration and published in Cuban journals. A
product without a stringent-authority approval cannot be sold in the markets that
pay, cannot be licensed to a partner who sells there, and cannot be priced at
anything like its therapeutic value. The COVID vaccines are the clearest case:
Cuba developed vaccines that worked, and could not sell them, because they did
not complete the international prequalification pathway.

\emph{No distribution partner.} Value in pharmaceuticals is captured through
regulatory affairs, distribution, salesforce and reimbursement negotiation ---
capabilities Cuba does not have and would take twenty years to build. The answer
is to partner, and Cuba has largely not partnered, in part because of the next
two reasons.

\emph{No capital.} Bringing one product through international trials costs more
than the whole sector's annual budget, and a defaulted sovereign cannot finance
it.

\emph{An ownership structure nobody will contract with.} A foreign
pharmaceutical company entering a licensing deal is buying a stream of
obligations enforceable over fifteen years. Its counterparty in Cuba is a state
holding company, in a jurisdiction with no independent courts, subject to a
statute that exposes anyone trafficking in confiscated property to litigation in
the United States, with a sovereign that froze foreign firms' accounts in 2009.
The science is not the risk.

\emph{And the regulator is the producer.} The Cuban medicines regulator, CECMED,
is a state body regulating products made by state bodies, within a single system
answerable to one government. Whatever the integrity of its individual staff ---
and its technical reputation is not bad --- it is structurally not independent,
and structural independence is precisely what a foreign regulator or a foreign
partner is looking for. This is the same conflict of interest
Chapter~\ref{ch:energy} identified in electricity, where the operator, owner,
planner and regulator are one entity.

\section{The design}

The strategy is to convert an existing scientific asset into an investable,
regulated, exporting sector without destroying the closed loop that made it work.
The order matters: regulator first, because nothing else is possible until a
foreign counterparty can rely on something.

\begin{design}{An independent medicines regulator}
\label{d:cecmed}
\begin{rules}
\rl{1}{Separation} The medicines regulator is constituted as an independent
authority with the tenure, budget and publication guarantees of
Design~\ref{d:courts}, institutionally separate from BioCubaFarma, from the
health ministry as owner, and from any producing entity.

\rl{2}{Fee-funded and published} Funded by application fees at published rates,
with all assessment reports, approvals, refusals and inspection findings
published, as the European and American agencies publish theirs. Publication is
what makes a regulator's judgments checkable by outsiders, which is the entire
asset being built.

\rl{3}{Target: stringent-authority status} An explicit, dated programme to reach
World Health Organization listed-authority standing, with the assessment
criteria, the gap analysis and the progress published annually. This is a
multi-year technical project with a definite endpoint, and it is the single
highest-return institutional investment available to Cuba in any sector.

\rl{4}{Reliance in the interim} Automatic recognition of approvals by stringent
foreign regulators for products Cuba imports, so that scarce assessor capacity is
spent on Cuban products seeking export rather than on re-reviewing medicines
already approved elsewhere. This also, immediately, shortens the queue for
imported medicines that are currently absent from pharmacies.

\rl{5}{Trials to international standard} All Cuban clinical trials conducted and
registered to international good-clinical-practice standards, in international
registries, with independent monitoring and ethics review. Cuban trial data is
currently discounted abroad for reasons of procedure rather than of substance,
and this clause is what makes forty years of accumulated science sellable.
\end{rules}
\end{design}

\begin{design}{Capital and partnership}
\label{d:biopartner}
\begin{rules}
\rl{1}{Product-level joint ventures} Partnership at the level of the individual
product --- one asset, one partner, defined territories, defined milestones ---
rather than at the level of the institution. This raises capital and buys
regulatory and distribution capability without selling the research base.

\rl{2}{Keep the platform, license the product} Cuba retains ownership of its
research institutions, its manufacturing and its underlying technology
platforms, and licenses rights to specific products in specific markets. The
comparators are Singapore and Ireland --- countries that built pharmaceutical
sectors by hosting and partnering rather than by originating --- and not
Switzerland, which is not an available model for a country of nine million with
no capital.

\rl{3}{Enforceable terms} Contracts under foreign law with international
arbitration, per Chapter~\ref{ch:capital}. A licensing agreement justiciable
only in Cuban courts is worth a fraction of the same agreement justiciable in
London or Singapore, and the discount is larger than anything Cuba could
negotiate on the royalty.

\rl{4}{Diaspora scientists} Cuban scientists abroad are the natural bridge to
foreign partners, regulators and capital, and Design~\ref{d:diaspora}.5 applies
directly. Several of the people best placed to take a Cuban product through an
FDA submission are Cubans who left.

\rl{5}{The domestic supply first} Nothing in this design permits export of a
medicine that is unavailable in Cuban pharmacies. Domestic supply is the first
call on production, stated as a rule because the revenue pressure will run the
other way.
\end{rules}
\end{design}

\section{The medical missions, reformed}

Cuba's largest export for over a decade was the labour of its health workers, and
Chapter~\ref{ch:venezuela} reported both halves of that: a real service really
delivered in places no other country would staff, financed by paying the workers
a fraction of the contract value and, in several destinations, holding their
passports.

The programme should be kept. It is valuable, it is genuinely good for the
receiving countries, and it is one of the few things Cuba sells that nobody else
sells. But the labour arrangement has to change, and the argument for changing it
is not only ethical. It is that the industry's binding constraint is now the
supply of willing doctors --- Cuban health workers are leaving the country
altogether, and a programme that treats them as conscripts loses them to
emigration, which yields Cuba nothing at all.

\begin{design}{Medical service export}
\label{d:missions}
\begin{rules}
\rl{1}{Individual contracts} Each health worker contracts individually, with the
full contract value disclosed to them and a published minimum share payable to
the worker. The state may retain a share --- there is a legitimate public claim on
the return to a publicly funded training --- and the share is published,
legislated, and the same for everyone.

\rl{2}{No control of documents or movement} Passports are held by their owners.
Restrictions on residence, association and relationships are abolished. The
penalty for leaving a mission is contractual, not a bar on returning to one's
own country, and Design~\ref{d:diaspora}.2 applies.

\rl{3}{Voluntary} Participation is voluntary, refusal carries no professional
consequence, and this is auditable.

\rl{4}{Compliant with labour standards} The programme is brought into conformity
with international labour conventions, and the conformity is externally verified
and published. This is what converts a programme under permanent international
criticism into an exportable, expandable and defensible service.

\rl{5}{Domestic staffing floor} A published minimum staffing level for the Cuban
health service, below which missions are not staffed. Per
Chapter~\ref{ch:premises}, people are the binding scarcity, and the domestic
system is the asset the whole book is written to protect.
\end{rules}
\end{design}

\section{Two further lines of business}

\emph{Clinical trials as an export.} Cuba has attributes that contract research
organisations pay for: a population with universal health coverage and complete
medical records, high follow-up and retention because people do not move or lose
their insurance, a physician-dense primary care network organised by defined
population, and a trial cost base far below North America or Europe. Under
Design~\ref{d:cecmed}.5 and an independent ethics framework, running phases of
international trials in Cuba is a genuine, capital-light, high-skill export that
also brings the domestic system access to investigational treatments. It requires
scrupulous consent standards --- a population that has limited access to medicines
is a population that can be induced to enrol --- and the ethics review has to be
independent of both the sponsor and the state, which is a stronger requirement
than most middle-income trial hosts apply.

\emph{Health tourism.} Cuba has treated foreign patients for decades and could
do far more of it, in cardiology, orthopaedics, ophthalmology, oncology,
neurological rehabilitation and addiction treatment, at prices well below the
United States and with outcomes that stand comparison. The constraint is
Chapter~\ref{ch:premises}'s again: every foreign patient consumes a Cuban
clinician's time. Design~\ref{d:delivery}.3 already binds --- designated
facilities, ring-fenced staff, revenue to the public system, and the programme
stops if the ring-fence cannot be held. This chapter adds nothing to that except
the observation that health tourism is the most valuable visitor per bed-night
available under Chapter~\ref{ch:tourism}'s metric, and the least environmentally
damaging.

\section{What is being protected}

A closing caution, because the reforms above are all in the direction of markets
and partnership and that direction has a known failure mode.

The Cuban biotechnology sector works because it is a closed loop with a public
customer, patient capital and no obligation to select targets by market size. Open
it to partnership incorrectly and it becomes a contract manufacturer for
somebody else's pipeline, with its research agenda set in New Jersey and its
public-health function abandoned --- at which point Cuba has sold its one
originating capability for a manufacturing margin.

The clauses that guard against that are Design~\ref{d:biopartner}.2, which keeps
the platform and licenses only products; Design~\ref{d:biopartner}.5, which puts
the domestic pharmacy first; and the public research budget, which
Chapter~\ref{ch:welfare}'s floor protects. They are the whole of the protection,
they are not self-enforcing, and a future government under fiscal pressure will
be offered a great deal of money to weaken them.
